\documentclass[11pt]{article}

\usepackage[T1]{fontenc}
\usepackage[margin=1in]{geometry}
\usepackage{amsmath,amssymb}
\usepackage{hyperref}

\title{Reliability Under Measurement Model Misspecification Is Target-Relative:\\
A Short Note on Omega Simulation Studies}
\author{}
\date{\today}

\newcommand{\Var}{\operatorname{Var}}
\newcommand{\Cov}{\operatorname{Cov}}
\newcommand{\T}{^{\mathsf T}}

\begin{document}
\maketitle

\begin{abstract}
Simulation studies of coefficient omega under measurement model
misspecification usually begin by choosing a population factor model, computing
the omega implied by that model, fitting alternative models to generated data,
and reporting bias relative to the population omega. This is a coherent
model-relative exercise. It is not, however, a model-free assessment of bias for
``the'' reliability of the composite. The reason is familiar from the literature
on equivalent structural equation models: the observed covariance matrix does
not, by itself, identify the substantive latent decomposition. Since omega as a
construct-reliability coefficient is the proportion of composite variance
attributed to a named target factor, equivalent decompositions can imply the
same item covariance matrix while assigning different reliabilities to the same
composite. Consequently, reliability under misspecification is difficult to pin
down unless the target construct, not only the observed covariance structure, is
declared. We use this point to reinterpret recent omega simulations as
conditional sensitivity analyses and suggest how future simulations can separate
estimator behavior from target uncertainty.
\end{abstract}

\noindent\textbf{Keywords:} coefficient omega, reliability, equivalent models,
measurement model misspecification, bifactor model

\section{The Problem}

Coefficient omega is attractive because it asks a substantively meaningful
question: how much of a composite score is attributable to the latent variable
the score is intended to measure? That question is also what makes omega
fragile. The answer depends on which latent variable is treated as the target
construct and which parts of the covariance structure are treated as nuisance,
method, local dependence, specific-factor variance, or error.

Bell, Chalmers, and Flora \cite{bell2024} give a useful recent example. They
simulate data from several population models, including a one-factor model with
correlated errors, bifactor models, and higher-order models. They then compare
omega estimates from fitted models with the omega implied by the population
model. Their central message is cautious: omega estimates are accurate when the
measurement model is correct, and can be biased when the model is misspecified.

The present note argues that there is a deeper target issue behind such
simulation results. Once a population model and target factor are stipulated,
the bias calculation is well-defined. But the word ``true'' is doing more work
than it may appear to do. It supplies not only a covariance matrix, but also a
substantive allocation of covariance to the target factor rather than to other
latent or residual sources. Under model misspecification, that allocation is
often exactly what is in question.

This is not a new problem in structural equation modeling. Equivalent models can
fit the same covariance matrix while representing different substantive
stories \cite{bollen1989,maccallum1993}. Fit alone cannot determine which story
is the data-generating one. The same difficulty becomes sharper for reliability
because the estimand itself can change across equivalent stories. Reliability is
not always a single property of a test; it can depend on what the score is meant
to measure \cite{borsboom2002,ellis2021,savalei2019}.

\section{Target-Factor Reliability}

Let \(y\) be the vector of item scores and let \(X=w\T y\) be the composite,
with \(w\) often equal to a vector of ones. Suppose a model says that a target
factor \(f\), scaled to have variance one, contributes loadings \(a\) to the
items:
\[
  y=a f+u,\qquad \Cov(f,u)=0,\qquad \Var(y)=\Sigma .
\]
The target-factor omega for \(X\) is
\begin{equation}\label{eq:omega}
  \omega_f
  =
  \frac{\Var(w\T a f)}{\Var(w\T y)}
  =
  \frac{(w\T a)^2}{w\T\Sigma w}.
\end{equation}
This definition uses the full composite variance in the denominator. If the
model contains correlated residuals, specific factors, or lower-order factors,
their contribution remains in \(w\T\Sigma w\). What changes from model to model
is the numerator: which part of the item covariance is assigned to the named
target factor?

This distinction is easy to miss. A simulation can label a population
reliability as .60 or .85, but that label is not a property of the covariance
matrix alone. It is a property of the covariance matrix together with a
particular latent interpretation.

\section{Equivalent Models Can Change the Target}

The issue can be seen with a simple construction. It is not meant as a proposed
measurement model. It is only a way to show what is, and is not, identified by
the covariance matrix. Let \(\Sigma\) be any positive definite item covariance
matrix. For any value \(r\) between zero and one, set
\[
  a_r=\sqrt{r}\frac{\Sigma w}{\sqrt{w\T\Sigma w}},
  \qquad
  \Psi_r=\Sigma-a_r a_r\T .
\]
Then
\[
  \Sigma=a_r a_r\T+\Psi_r
\]
is a one-factor model with correlated residuals. The residual covariance
\(\Psi_r\) is nonnegative definite by Cauchy's inequality, and
\((w\T a_r)^2=r\,w\T\Sigma w\). Equation~\eqref{eq:omega} therefore gives
\(\omega_{f_r}=r\) for the composite \(X=w\T y\).

Thus, if correlated residuals are allowed without further restriction, the same
observed covariance matrix can be paired with many different target-factor
reliabilities. An omega estimator that converges to a fixed function of
\(\Sigma\) can therefore be positively biased, unbiased, or negatively biased
depending only on which equivalent latent decomposition is called the population
model.

This is the sense in which reliability under misspecification is hard to pin
down. The issue is not only that the fitted model may be wrong. It is that the
target against which the fitted model is judged may not be identified without
extra theory.

\section{Re-reading Omega Misspecification Simulations}

The point above does not make Bell et al.'s \cite{bell2024} simulations
incorrect. Their population models define clear conditional targets. For a
one-factor correlated-error population, the target is the part of composite
variance due to the one common factor. For a bifactor population, the target is
the part due to the general factor. For a higher-order population, the target is
the part due to the higher-order factor through its indirect effects.

The limitation is that these are target-relative claims. Consider a bifactor
population,
\[
  \Sigma=gg\T+SS\T+\Theta ,
\]
where \(g\) is the general-factor loading vector and \(S\) contains specific
factor loadings. The bifactor omega treats \(gg\T\) as construct-relevant
general-factor variance and treats \(SS\T+\Theta\) as non-target variance for
the total score. Marginally, however, \(SS\T+\Theta\) is a structured residual
covariance matrix around the general factor. Another one-factor
correlated-residual decomposition can imply the same \(\Sigma\) but assign a
different amount of composite variance to its common factor. The observed item
covariance matrix cannot decide which target is substantively correct.

The higher-order case illustrates the important qualification. A higher-order
model and a restricted bifactor model can be related by the Schmid-Leiman
transformation \cite{schmid1957,yung1999}. When that transformation preserves
the vector of general-factor contributions to the items, omega is preserved too.
Equivalent models are not automatically a problem for reliability. They become a
problem when the equivalent representations change the target contribution
vector in Equation~\eqref{eq:omega}.

Thus, simulations of misspecified omega estimates are best read as answering a
conditional question:
\[
  \text{If this latent target story is true, how does this fitted omega behave?}
\]
They do not answer the unconditional question:
\[
  \text{How biased is omega for the reliability encoded in the covariance matrix?}
\]
There is no such covariance-only target without further restrictions.

\section{Implications}

The practical implication is not that omega should be abandoned, or that
misspecification simulations are useless. The implication is that simulations
should separate two problems that are often blended together.

First, there is an estimator problem. Given a target model, a target factor, and
a population covariance matrix, what does a fitted omega estimate in finite
samples or asymptotically? Bell et al. address this problem.

Second, there is a target problem. Given the same observed covariance matrix,
which part of the composite variance should count as construct-relevant
reliability? This problem is not solved by Monte Carlo replication, larger
sample size, or better fit indices. It requires theory, design information,
external validation, or explicit sensitivity analysis over plausible equivalent
or near-equivalent decompositions.

Future simulation studies can make this distinction explicit by reporting their
results as conditional on a declared target factor, by noting whether the target
is invariant across relevant equivalent models, and by including sensitivity
conditions that move covariance between general factors, specific factors, and
correlated residuals while holding the item covariance matrix fixed or nearly
fixed. Such designs would not merely ask whether an estimator is biased under
misspecification. They would ask whether the proposed reliability target itself
is stable under the kinds of model uncertainty that applied researchers face.

\section{Conclusion}

The established lesson from equivalent-model work in SEM is that the same
covariance matrix can support different substantive explanations. For
model-based reliability coefficients, the consequence is stronger: different
explanations can define different reliabilities for the same composite. Bias
under misspecification is therefore not a model-free property of coefficient
omega. It is a property of an estimator relative to a declared latent target.
That framing preserves the value of recent omega simulations while clarifying
their scope. They are conditional sensitivity studies, not definitive evidence
about bias for a uniquely identified reliability of the observed composite.

\begin{thebibliography}{9}

\bibitem{bell2024}
Bell, S. M., Chalmers, R. P., \& Flora, D. B. (2024).
The impact of measurement model misspecification on coefficient omega estimates
of composite reliability.
\emph{Educational and Psychological Measurement}, 84(1), 5--39.
\href{https://doi.org/10.1177/00131644231155804}{doi:10.1177/00131644231155804}.

\bibitem{bollen1989}
Bollen, K. A. (1989).
\emph{Structural equations with latent variables}.
Wiley.

\bibitem{borsboom2002}
Borsboom, D., \& Mellenbergh, G. J. (2002).
True scores, latent variables, and constructs: A comment on Schmidt and Hunter.
\emph{Intelligence}, 30(6), 505--514.
\href{https://doi.org/10.1016/S0160-2896(02)00082-X}{doi:10.1016/S0160-2896(02)00082-X}.

\bibitem{ellis2021}
Ellis, J. L. (2021).
A test can have multiple reliabilities.
\emph{Psychometrika}, 86(4), 869--876.
\href{https://doi.org/10.1007/s11336-021-09800-2}{doi:10.1007/s11336-021-09800-2}.

\bibitem{maccallum1993}
MacCallum, R. C., Wegener, D. T., Uchino, B. N., \& Fabrigar, L. R. (1993).
The problem of equivalent models in applications of covariance structure
analysis.
\emph{Psychological Bulletin}, 114(1), 185--199.
\href{https://doi.org/10.1037/0033-2909.114.1.185}{doi:10.1037/0033-2909.114.1.185}.

\bibitem{savalei2019}
Savalei, V., \& Reise, S. P. (2019).
Don't forget the model in your model-based reliability coefficients: A reply to
McNeish (2018).
\emph{Collabra: Psychology}, 5(1), Article 36.
\href{https://doi.org/10.1525/collabra.247}{doi:10.1525/collabra.247}.

\bibitem{schmid1957}
Schmid, J., \& Leiman, J. M. (1957).
The development of hierarchical factor solutions.
\emph{Psychometrika}, 22(1), 83--90.
\href{https://doi.org/10.1007/BF02289209}{doi:10.1007/BF02289209}.

\bibitem{yung1999}
Yung, Y. F., Thissen, D., \& McLeod, L. D. (1999).
On the relationship between the higher-order factor model and the hierarchical
factor model.
\emph{Psychometrika}, 64(2), 113--128.
\href{https://doi.org/10.1007/BF02294531}{doi:10.1007/BF02294531}.

\end{thebibliography}

% Development note: before submission, add a fuller review of equivalent-model
% results specific to CFA, local dependence, bifactor models, and reliability
% under multidimensional or method-factor interpretations.

\end{document}
